% ============================================================
%  PREAMBULO.TEX
%  Configuración maestra de la Monografía UMSS — Dirección de Posgrado
%  Motor requerido: LuaLaTeX (recomendado) o XeLaTeX
%  ------------------------------------------------------------
%  Todo lo que define el "look" del documento vive AQUI.
%  Las secciones de contenido (secciones/*.tex) NO deben tocar
%  formato: solo texto. Así, un cambio de estilo se hace UNA vez.
% ============================================================

% ---------- 1. CLASE BASE ----------------------------------
% Se usa memoir: es la clase con mejor soporte nativo para
% frontmatter/mainmatter (numeración romana -> arábiga),
% control fino de índices (TOC/LOF/LOT) y estilos de "caption",
% todo sin necesitar media docena de paquetes que compitan entre sí
% (evita choques típicos de report+tocloft+titlesec+fancyhdr).
\documentclass[11pt,oneside,letterpaper]{memoir}
% Cuerpo en 11,5 pt: memoir solo ofrece tamaños enteros, así que se
% parte de 11pt (que fija \small, \footnotesize, etc.) y se redefine
% únicamente \normalsize a 11,5/13,8 pt, con los mismos espacios de
% ecuaciones que size11.
\makeatletter
\renewcommand*{\normalsize}{%
  \@setfontsize\normalsize{11.5}{13.8}%
  \abovedisplayskip 11\p@ \@plus3\p@ \@minus6\p@
  \abovedisplayshortskip \z@ \@plus3\p@
  \belowdisplayshortskip 6.5\p@ \@plus3.5\p@ \@minus3\p@
  \belowdisplayskip \abovedisplayskip
  \let\@listi\@listI}
\normalsize
\makeatother

% ---------- 2. IDIOMA Y CODIFICACIÓN ------------------------
% Con LuaLaTeX/XeLaTeX el motor ya trabaja en UTF-8 nativo;
% NO se cargan inputenc/fontenc (son innecesarios y darían
% advertencias de obsolescencia en un motor moderno).
\usepackage{polyglossia}
\setmainlanguage{spanish}
% Alternativa equivalente si se prefiere babel:
% \usepackage[spanish,es-noindentfirst]{babel}

% ---------- 3. MARGENES (Instrucción 3, formato general) ----
% Izq 3cm / Der 2.5cm / Sup 2.5cm / Inf 2.5cm, carta (letter),
% medidos sobre el BLOQUE DE TEXTO. El encabezado y el folio
% viven DENTRO de los márgenes (guía 3.6): no descuentan del
% cuerpo. Se usa el sistema de maquetación nativo de memoir (no
% geometry) para que ninguna opción "includehead/includefoot"
% vuelva a comerse 1,3 cm de cada margen.
\setstocksize{11in}{8.5in}
\settrimmedsize{\stockheight}{\stockwidth}{*}
\setlrmarginsandblock{3cm}{2.5cm}{*}
\setulmarginsandblock{2.5cm}{2.5cm}{*}
\setheadfoot{14pt}{1.25cm}     % alto del encabezado; distancia última línea -> folio
\setheaderspaces{*}{0.75cm}{*} % separación encabezado -> primera línea del texto
\checkandfixthelayout[fixed]   % [fixed]: no reajusta el bloque a múltiplos de línea

% ---------- 4. TIPOGRAFIA (fontspec) ------------------------
% pdflatex NO puede cargar fuentes TrueType/OpenType del sistema
% (Arial, Garamond) de forma nativa: solo soporta fuentes Type1
% de 8 bits, obliga a "fingir" Arial con métricas de sustitución
% manuales, y no tiene soporte real de OpenType (kerning, small
% caps, ligaduras contextuales) ni de Unicode de entrada. Por eso
% se exige LuaLaTeX/XeLaTeX + fontspec, que cargan la fuente real
% instalada en el sistema (o su clon métricamente compatible).
\usepackage{fontspec}
\usepackage{microtype}   % micro-tipografía: protrusion, expansion, kerning fino

% --- Fuente del cuerpo: Garamond 11,5 pt ---------------------
% El cuerpo, las tablas, los rótulos y las notas van en Garamond.
% Se carga EB Garamond instalada en el sistema (paquete AUR
% ebgaramond-otf) POR NOMBRE DE ARCHIVO, no por nombre de familia:
% así no depende del índice de luaotfload, que era la causa de los
% fallos al instalar fuentes nuevas. Cascada de respaldo: Garamond
% real (Windows/Overleaf) -> EB Garamond por archivo -> por nombre.
% \opcionesCuerpo guarda la misma selección para la matemática.
% \mathRecta/\mathItalica: la misma fuente, redonda y cursiva, para
% la matemática (\setmathfont no acepta UprightFont/ItalicFont).
\IfFontExistsTF{Garamond}
  {\def\fuenteCuerpo{Garamond}\def\opcionesCuerpo{}%
   \def\mathRecta{Garamond}\def\mathItalica{Garamond Italic}\def\opcionesMath{}}
  {\IfFileExists{/usr/share/fonts/EBGaramond12-otf/EBGaramond-Regular.otf}
    {\def\fuenteCuerpo{EBGaramond}%
     \def\opcionesCuerpo{Path=/usr/share/fonts/EBGaramond12-otf/,
       Extension=.otf, UprightFont=*-Regular, BoldFont=*-Bold,
       ItalicFont=*-Italic, BoldItalicFont=*-BoldItalic}%
     \def\mathRecta{EBGaramond-Regular}\def\mathItalica{EBGaramond-Italic}%
     \def\opcionesMath{Path=/usr/share/fonts/EBGaramond12-otf/, Extension=.otf,}}
    {\def\fuenteCuerpo{EB Garamond}\def\opcionesCuerpo{}%
     \def\mathRecta{EB Garamond}\def\mathItalica{EB Garamond Italic}\def\opcionesMath{}}}
\expandafter\setmainfont\expandafter{\fuenteCuerpo}[Numbers=Lining,\opcionesCuerpo]
% Alias heredado: el Resumen y el glosario usan \garamondfont.
\newcommand{\garamondfont}{\normalfont}

% --- Fuente de títulos: Arimo (clon métrico de Arial) -------
% Títulos, índices, encabezado/pie y carátula en sans-serif, como
% pide la guía. Arimo tiene los mismos anchos que Arial. Se usa la
% versión Nerd Font del sistema (ttf-arimo-nerd), también por nombre
% de archivo; cascada: Arimo Nerd Font -> Arimo -> Arial -> Liberation Sans.
\IfFileExists{/usr/share/fonts/TTF/ArimoNerdFont-Regular.ttf}
  {\newfontfamily\arialfont{ArimoNerdFont}[
     Path           = /usr/share/fonts/TTF/,
     Extension      = .ttf,
     UprightFont    = *-Regular,
     BoldFont       = *-Bold,
     ItalicFont     = *-Italic,
     BoldItalicFont = *-BoldItalic,
     Ligatures      = TeX]}
  {\IfFontExistsTF{Arimo}
    {\newfontfamily\arialfont{Arimo}[Ligatures=TeX]}
    {\IfFontExistsTF{Arial}
      {\newfontfamily\arialfont{Arial}[Ligatures=TeX]}
      {\newfontfamily\arialfont{Liberation Sans}[Ligatures=TeX]}}}

% --- Monoespaciada (variables, código): IBM Plex Mono -------
% Libre (OFL) y empaquetada en fuentes/ para que el PDF sea igual
% en local, Overleaf y CI. Se escala a la altura de x de Garamond
% para que "origin_latitude" no parezca más grande que el texto.
\setmonofont{IBMPlexMono}[
  Path           = fuentes/,
  Extension      = .otf,
  UprightFont    = *-Regular,
  BoldFont       = *-Bold,
  ItalicFont     = *-Italic,
  BoldItalicFont = *-BoldItalic,
  Scale          = MatchLowercase]

% --- Matemática en Garamond ---------------------------------
% Las expresiones del texto (p = 0,033; ρ; D²; 10⁻⁹) toman cifras
% y letras latinas y griegas de la fuente del cuerpo; los operadores
% y símbolos (∈, ≤, →, Σ) salen de Latin Modern Math, cuyo trazo
% serif combina con Garamond.
\usepackage{amsmath}
\usepackage[math-style=TeX]{unicode-math}
\setmathfont{Latin Modern Math}
\edef\x{\noexpand\setmathfont{\mathRecta}[\opcionesMath
  range={up/{num,latin,Latin,greek,Greek}}]}\x
\edef\x{\noexpand\setmathfont{\mathItalica}[\opcionesMath
  range={it/{latin,Latin,greek,Greek}->up}]}\x
\usepackage{icomma}  % "0,5" en modo matemático sin espacio tras la coma

% ---------- 5. INTERLINEADO ---------------------------------
% Cuerpo en Garamond 11,5 pt con interlineado múltiple 1,2, el mismo
% del Resumen: \normalsize tiene \baselineskip de 13,8 pt (1,2 × 11,5)
% y \setSpacing{1.2} lo lleva a unos 16,6 pt entre líneas.
% \setSpacing (sin asterisco) deja en sencillo flotantes (tablas,
% figuras, captions, "Nota.") y notas al pie.
% NOTA TÉCNICA: \setSpacing solo toma efecto al final de un cambio
% de tamaño de fuente (reinvoca \@currsize), y en el preámbulo esa
% macro todavía no está ligada a \normalsize. Por eso el interlineado
% del cuerpo se activa en \AtBeginDocument, ya con la fuente cargada.
\newcommand{\interlineadoCuerpo}{\setSpacing{1.2}}  % cuerpo: 1,2
\AtBeginDocument{\interlineadoCuerpo}
\newcommand{\interlineadoResumen}{\setSpacing{1.2}} % Resumen: 1.2 (Sección preliminar)
\newcommand{\interlineadoSimple}{\setSpacing{1}}    % carátula, bibliografía interna, notas

% ---------- 6. NUMERACIÓN, ENCABEZADO Y PIE DE PÁGINA -------
% memoir define \frontmatter (activa romanos i,ii,iii...) y
% \mainmatter (reinicia en arábigo 1,2,3...) automáticamente:
% así se cumple 3.1 sin tocar contadores a mano.
% La POSICIÓN del folio (guía §3.1: abajo y centrada) la fija
% el interruptor "folioDerecha" del panel de control.
% Tres estilos:
%   preliminar  carátula a glosario: solo el folio.
%   cuerpo      desde \mainmatter: encabezado corrido (capítulo a la
%               izquierda, sección a la derecha, filete fino) y pie con
%               filete, autor, folio e institución/año.
%   plain       primera hoja de cada capítulo: sin encabezado (el título
%               del capítulo ya está en la página), mismo pie que "cuerpo".
% El encabezado y el pie van en Arimo 8-9 pt y gris, para que no
% compitan con el texto; viven dentro de los márgenes (sección 3).
\usepackage{fancyhdr}
\usepackage[breakwords]{truncate}  % recorta títulos largos con "…"
\newcommand{\estiloMarcas}{\arialfont\footnotesize\color{gray!80!black}}
\newcommand{\estiloPie}{\arialfont\fontsize{8}{10}\selectfont\color{gray!80!black}}
\newcommand{\folio}{{\arialfont\normalsize\thepage}}
% Marcas: "7. Desarrollo" y "7.5 Evaluación y Resultados", tal como
% se escribieron (sin forzar mayúsculas). Al abrir un capítulo la
% marca de sección se vacía hasta la primera \section.
\newcommand{\pieComun}{%
  \fancyfoot[L]{\estiloPie\nombrePostulante}%
  \fancyfoot[C]{\folio}%
  \fancyfoot[R]{\estiloPie UMSS\,·\,\anioMonografia}%
  \iffolioDerecha
    \fancyfoot[C]{}%
    \fancyfoot[R]{\estiloPie UMSS\,·\,\anioMonografia\quad\folio}%
  \fi
  \renewcommand{\footrulewidth}{0.3pt}}
\fancypagestyle{cuerpo}{%
  \fancyhf{}%
  \fancyhead[L]{\estiloMarcas\truncate{0.49\textwidth}{\leftmark}}%
  \fancyhead[R]{\estiloMarcas\truncate{0.49\textwidth}{\rightmark}}%
  \renewcommand{\headrulewidth}{0.4pt}%
  \pieComun}
\fancypagestyle{plain}{%
  \fancyhf{}%
  \renewcommand{\headrulewidth}{0pt}%
  \pieComun}
\fancypagestyle{preliminar}{%
  \fancyhf{}%
  \fancyfoot[C]{\folio}%
  \iffolioDerecha\fancyfoot[C]{}\fancyfoot[R]{\folio}\fi
  \renewcommand{\headrulewidth}{0pt}%
  \renewcommand{\footrulewidth}{0pt}}
% En los preliminares también las aperturas (Resumen, Índice...) van
% solo con folio; main.tex cambia a "cuerpo"/"plain" tras \mainmatter
% con \estiloCuerpo.
\pagestyle{preliminar}
\aliaspagestyle{chapter}{preliminar}
\aliaspagestyle{part}{preliminar}
\newcommand{\estiloCuerpo}{%
  \pagestyle{cuerpo}%
  \aliaspagestyle{chapter}{plain}%
  \aliaspagestyle{part}{plain}}
\makeatletter
\AtBeginDocument{%
  \renewcommand{\chaptermark}[1]{%
    \if@mainmatter\markboth{\thechapter.\ #1}{}\else\markboth{#1}{}\fi}%
  \renewcommand{\sectionmark}[1]{\markright{\thesection\ #1}}}
\makeatother
% Copia del número de la última hoja preliminar: permite que el
% interruptor "numeracionContinua" reanude la serie arábica
% donde terminó el índice de figuras, en vez de reiniciarla.
\newcounter{pagPreliminar}
\newcommand{\guardarPagPreliminar}{\setcounter{pagPreliminar}{\value{page}}}

% ---------- 7. TITULOS DE SECCIONES (Instrucción 3) ---------
% "Títulos principales en Arial 14 y negrilla": se usa \chapter
% para las 9 secciones principales del Anexo A (fuerza salto de
% página automático, como exige la guía) y \section para el
% siguiente nivel (2.1, 5.1, 7.1, etc.).
\makeatletter
\setlength{\beforechapskip}{0pt}
\setlength{\midchapskip}{6pt}
\setlength{\afterchapskip}{30pt}
% NOTA TÉCNICA: memoir reutiliza \printchaptertitle tanto para
% capítulos numerados como para los títulos SIN número del
% índice general / índice de figuras / índice de tablas (los
% genera internamente como si fueran "capítulos"). Por eso el
% número ("1.", "2.", ...) se resuelve en \printchapternum (que
% memoir SOLO invoca para capítulos numerados reales) y NO
% dentro de \printchaptertitle: así "ÍNDICE"/"ÍNDICE DE FIGURAS"
% nunca arrastran un "0." espurio.
\makechapterstyle{umss}{
  \renewcommand{\chaptitlefont}{\arialfont\fontsize{14}{16}\selectfont\bfseries}
  \renewcommand{\chapterheadstart}{\vspace*{\beforechapskip}}
  % Camino NUMERADO (\chapter{...} normal):
  \renewcommand{\printchaptername}{\chaptitlefont\centering}% activa fuente+centrado (sin texto)
  \renewcommand{\chapternamenum}{}
  \renewcommand{\printchapternum}{\thechapter.\space}% "N. "
  \renewcommand{\afterchapternum}{}% SIN salto de línea: numero y título van juntos
  % Camino SIN NÚMERO (índice general, índice de figuras/tablas):
  \renewcommand{\printchapternonum}{\chaptitlefont\centering}% misma fuente+centrado, sin texto
  % Título (para ambos caminos, ya con la fuente activa):
  \renewcommand{\printchaptertitle}[1]{\MakeUppercase{##1}\par}
  \renewcommand{\afterchaptertitle}{\par\nobreak\vskip\afterchapskip}
}
\chapterstyle{umss}
\makeatother

% Subtítulos (\section, ej. "2.1 Descripción del Problema"):
% mismo criterio tipográfico que el título principal pero un
% cuerpo menor y alineado a la izquierda (ajustable aquí mismo).
\setsecheadstyle{\arialfont\fontsize{13}{15}\selectfont\bfseries}
\setbeforesecskip{-2.5ex plus -1ex minus -.2ex}
\setaftersecskip{1.3ex plus .2ex}
\setsubsecheadstyle{\arialfont\fontsize{12}{14}\selectfont\bfseries\itshape}

% Cuerpo general en Garamond 11,5: opción de clase [11pt] más la
% redefinición de \normalsize (sección 1) e interlineado 1,2
% (sección 5). Los demás tamaños (\small, \footnotesize...) son
% los de memoir 11pt.

% ---------- 8. INDICE GENERAL / LOF / LOT (3.4) --------------
% Reglas exigidas: sin negrillas, doble espacio entre entradas
% de nivel principal, puntos de relleno hasta el número de
% página, palabra "Página" sobre la columna de números,
% secciones principales sin sangría, subtítulos sangrados.
% IMPORTANTE: polyglossia (gloss-spanish.ldf) reasigna
% \contentsname -> "Índice general", \listtablename -> "Índice de
% cuadros" y \tablename -> "Cuadro" AL INICIAR EL DOCUMENTO, es
% decir DESPUÉS de este preámbulo. Por eso un \renewcommand suelto
% aquí se perdería silenciosamente, y ni siquiera \AtBeginDocument
% basta (polyglossia activa el idioma en una fase posterior).
% Se usa el gancho moderno "begindocument/end" del núcleo LaTeX,
% que se ejecuta al FINAL del arranque del documento: así estos
% nombres son los últimos en aplicarse y siempre ganan.
\AddToHook{begindocument/end}{%
  \renewcommand{\contentsname}{ÍNDICE}%
  \renewcommand{\listfigurename}{ÍNDICE DE FIGURAS}%
  \renewcommand{\listtablename}{ÍNDICE DE TABLAS}%
  % La guía (3.2) rotula "Tabla 1" y "Figura 1", no el "Cuadro"
  % que impone el español de polyglossia:
  \renewcommand{\tablename}{Tabla}%
  \renewcommand{\figurename}{Figura}%
}

% Profundidad de TRES niveles en la numeración y en el índice
% (capítulo · sección · subsección, ej. 7 · 7.2 · 7.2.1), tal
% como exige 2.4. \settocdepth también controla hasta qué nivel
% llegan los marcadores del PDF vía bookmarksdepth (más abajo).
\maxsecnumdepth{subsection}
\settocdepth{subsection}

% Nivel "chapter" (secciones principales 1..9): memoir ya trae
% ganchos "cft" propios para este nivel -> se renuevan directo.
% "Mismo formato de títulos que el texto" (2.4): los capítulos van en
% mayúsculas en el cuerpo (vía \printchaptertitle) y TAMBIÉN en el
% índice. Memoir invoca \cftchapterfont así: {\cftchapterfont {#1}} —
% por eso basta con que \cftchapterfont termine en \MakeUppercase (sin
% argumento propio): TeX toma el grupo {#1} que le sigue como su
% argumento. #1 incluye el numberline, pero \MakeUppercase no afecta a
% los dígitos, así que "7." no se ve alterado.
% Columna del folio en índices: memoir reserva 1.55em, que no alcanza
% para folios de tres cifras en Arial 12 (desborde de 1,4 pt por línea
% desde la página 100). Se amplía el folio y, en la misma medida, el
% margen derecho del texto de cada entrada.
\setpnumwidth{2.2em}
\setrmarg{3.2em}
\renewcommand{\cftchapterfont}{\arialfont\normalsize\MakeUppercase} % SIN negrilla (regla explícita)
\renewcommand{\cftchapterpagefont}{\arialfont\normalsize}
\cftsetindents{chapter}{0em}{1.5em}   % secciones principales SIN sangría
\renewcommand{\cftchapterleader}{\normalfont\cftdotfill{\cftdotsep}} % puntos, sin negrilla
\renewcommand{\cftchapterafterpnum}{}  % sin espacio tras el folio
\renewcommand{\cftchapterpresnum}{}
\renewcommand{\cftchapteraftersnum}{.}  % "1. Introducción" en el índice
\setlength{\cftbeforechapterskip}{14pt}  % doble espacio ANTES de cada capítulo (≈1 línea sencilla)

% Niveles "section"/"subsection"/"figure"/"table": memoir NO
% define ganchos "cft*font" para estos (solo para chapter/part/
% book); en vez de forzar el paquete tocloft encima de memoir
% (choca: ambos declaran \cftchapterfont, etc., con \newcommand
% no condicional) se redefine directamente el mecanismo clásico
% de LaTeX \l@<nivel>, que es el que memoir usa por debajo para
% estos niveles. Así se logra el mismo resultado (fuente Arial,
% sin negrilla, sangría, puntos de relleno) sin ningún paquete
% adicional ni conflicto de nombres.
% (\@dottedtocline es un comando interno con "@" en el nombre:
%  se necesita \makeatletter/\makeatother alrededor, aquí sí.)
\makeatletter
\renewcommand*{\l@section}[2]{%
  {\arialfont\normalsize\@dottedtocline{1}{1.5em}{2.3em}{#1}{#2}}}
\renewcommand*{\l@subsection}[2]{%
  {\arialfont\normalsize\@dottedtocline{2}{3em}{3em}{#1}{#2}}}
% Las entradas de \listoffigures/\listoftables solo traen el
% NÚMERO (\numberline{N}), no la palabra "Figura"/"Tabla": la
% guía (2.4) exige el prefijo también en el índice, no solo en el
% rótulo de la figura/tabla misma. Se redefine \numberline
% LOCALMENTE (dentro del grupo que abre cada \l@figure/\l@table)
% para anteponerlo; como el argumento #1 solo se expande cuando
% \@dottedtocline lo compone, la redefinición local sí surte
% efecto pese a estar "después" en el código.
\renewcommand*{\l@figure}[2]{%
  {\arialfont\normalsize\def\numberline##1{\figurename~##1\hspace{0.4em}}%
   \@dottedtocline{1}{0em}{4.2em}{#1}{#2}}}
\renewcommand*{\l@table}[2]{%
  {\arialfont\normalsize\def\numberline##1{\tablename~##1\hspace{0.4em}}%
   \@dottedtocline{1}{0em}{4.2em}{#1}{#2}}}
\makeatother

% Encabezado "Página" alineado a la derecha, sobre la columna
% de folios, tal como exige 3.4. Se invoca antes de cada índice.
% Las entradas del índice van a interlineado sencillo (guía 3.4),
% restaurando el doble al terminar.
\newcommand{\encabezadoIndice}{%
  \noindent\hfill{\arialfont\normalsize\bfseries Página}\par\vspace{2pt}%
  \interlineadoSimple%
}
\newcommand{\cierreFuenterIndice}{%
  \par\vspace{0pt}%  cierre silencioso del grupo de sencillo
  \interlineadoCuerpo%  restaura el del cuerpo
}
% memoir ofrece ganchos que se ejecutan justo DESPUÉS del título
% del índice y ANTES/DESPUÉS de las entradas: es el punto exacto
% donde debe aparecer la palabra "Página" y donde restauramos el doble.
% Así se aplica solo, en los tres índices, sin escribir nada en main.tex.
\renewcommand{\cfttocbeforelisthook}{\encabezadoIndice}
\renewcommand{\cfttocafterlisthook}{\cierreFuenterIndice}
\renewcommand{\cftlofbeforelisthook}{\encabezadoIndice}
\renewcommand{\cftlofafterlisthook}{\cierreFuenterIndice}
\renewcommand{\cftlotbeforelisthook}{\encabezadoIndice}
\renewcommand{\cftlotafterlisthook}{\cierreFuenterIndice}

% ---------- 9. TABLAS Y FIGURAS (3.2) ------------------------
% Regla visual exigida: "Tabla N" en negrita centrado, debajo el
% título en cursiva centrado (caption ARRIBA del contenido),
% cuerpo de la tabla con SOLO líneas horizontales (booktabs),
% y una línea "Fuente: ..." en cursiva debajo. Como el orden
% físico de \caption y \includegraphics/tabular dentro del
% entorno determina si el título queda arriba o abajo, basta con
% escribir \caption{} ANTES del contenido (ver plantillas de
% ejemplo en cada sección).
% Flotantes anclados a su sección (guía 2.5): \FloatBarrier antes de
% cada \section para que las figuras no escapen a secciones posteriores.
\usepackage{placeins}
\usepackage{etoolbox}
\AtBeginEnvironment{section}{\FloatBarrier}
% Espacio vertical y parámetros de flotación ajustados para evitar
% páginas parcialmente vacías (aumenta \textfraction, reduce flotantes sueltos).
\setlength{\intextsep}{18pt plus 4pt minus 4pt}
\setlength{\textfloatsep}{20pt plus 4pt minus 4pt}
\setlength{\floatsep}{16pt plus 4pt minus 2pt}
\renewcommand\topfraction{.85}
\renewcommand\bottomfraction{.7}
\renewcommand\textfraction{.15}
\renewcommand\floatpagefraction{.7}
% Nota: memoir no usa \fps@figure/\fps@table; el especificador [htbp]
% se controla en el entorno directamente. La \FloatBarrier de placeins
% es suficiente para anclar flotantes a su sección.
\usepackage{booktabs}       % líneas horizontales limpias, sin verticales
\usepackage{graphicx}
\usepackage{caption}
% Especificador [H]: el flotante queda exactamente donde se escribe, entre dos párrafos.
% Con [htbp] LaTeX puede retrasar una figura al tope de la página siguiente, por
% encima de las líneas con que continúa el párrafo anterior.
\usepackage{float}
% Rótulo "Tabla N"/"Figura N" en negrita + título en cursiva, EN UNA
% SOLA LÍNEA (2.5: "rótulo arriba, en una línea, centrado"), separados
% por un espacio (antes iban en líneas distintas con labelsep=newline).
% Espaciado ajustado: figura+caption+nota juntos, sin salto de página.
\captionsetup{
  font=normalfont,
  labelfont=bf,
  textfont=it,
  justification=centering,
  singlelinecheck=false,
  labelsep=space,
  skip=4pt,
  aboveskip=4pt
}
\captionsetup[figure]{
  aboveskip=4pt
}
% La guía numera de forma CONSECUTIVA en todo el documento
% ("Tabla 1", "Tabla 2", ... "Figura 1", "Figura 2"), no por
% capítulo ("Tabla 1.1"). En memoir los contadores table/figure
% cuelgan de "chapter", así que no basta con redefinir \thetable:
% hay que DESACOPLAR el contador, o seguiría reiniciándose en
% cada sección principal (dando dos "Tabla 1" distintas).
% TRAMPA IMPORTANTE DE MEMOIR: el propio comando \mainmatter
% ejecuta internamente \counterwithin{table}{chapter} y
% \counterwithin{figure}{chapter}, es decir, RE-ACOPLA los
% contadores al capítulo en tiempo de ejecución y deshace lo que
% se haya configurado en el preámbulo (por eso aparecía
% "Tabla 1.1"). La solución estable es vaciar ese gancho interno
% (\@memmain@floats) para que \mainmatter ya no re-acople nada.
\counterwithout{table}{chapter}
\counterwithout{figure}{chapter}
\makeatletter
\renewcommand\@memmain@floats{}   % \mainmatter ya no re-acopla
\renewcommand\@memfront@floats{}  % \frontmatter tampoco toca los contadores
\makeatother
\renewcommand{\thetable}{\arabic{table}}
\renewcommand{\thefigure}{\arabic{figure}}

% MARCADOR DE POSICIÓN SEGURO PARA TABLAS
% Dentro de un "tabular", si una celda comienza con un corchete
% justo después de un salto de fila (\\), LaTeX interpreta ese
% corchete como el argumento OPCIONAL de \\ (un espaciado
% vertical) e intenta leer su contenido como una longitud: el
% documento queda esperando indefinidamente. Por eso los textos
% de relleno dentro de tablas se escriben con \ph{...}, que
% imprime el corchete de forma segura.
\newcommand{\ph}[1]{[#1]}

% "Nota." bajo tabla/figura (2.5): una sola línea, sencillo, 10 pt,
% con "Nota." en cursiva y el resto en texto normal -- reemplaza al
% antiguo rótulo "Fuente:" conservando el mismo nombre de comando
% (\fuente) para no tocar ningún archivo de secciones/.
% \nopagebreak mantiene figura+caption+nota juntos.
\newcommand{\fuente}[1]{%
  \nopagebreak
  \par\vspace{2pt}
  {\interlineadoSimple\fontsize{10}{10}\selectfont\textit{Nota.} #1\par}
}

% ---------- 10. CITAS Y REFERENCIAS (3.5 / 3.6) --------------
% biblatex + biber, con un único referencias.bib (Fase 4). La guía
% pide formato "americano" (Autor, año) en el texto -- lo que da
% \parencite de biblatex sin más -- y una lista final alfabética,
% sin numerar, con sangría francesa, apellido e iniciales y título
% en cursiva: es, en sustancia, el estilo APA, así que se usa
% biblatex-apa directamente en vez de un .bbx propio.
% El interruptor \ifbibNumerada (configuracion.tex, P3 opcional de
% 3.6 "número y autor en negrilla, apellido en MAYÚSCULAS") NO tiene
% un .bbx institucional propio todavía: alternar a "numeric" da una
% lista numerada correcta pero sin ese detalle tipográfico exacto;
% se deja así documentado como pendiente (ver informe final).
\makeatletter
\ifbibNumerada
  \PassOptionsToPackage{backend=biber,style=numeric,sorting=none}{biblatex}
\else
  \PassOptionsToPackage{backend=biber,style=apa,sorting=nyt}{biblatex}
\fi
\makeatother
\usepackage{csquotes}   % biblatex lo recomienda para comillas sensibles al idioma
\usepackage{biblatex}
\addbibresource{referencias.bib}
% Los \cita{Autor}{Año} artesanales se sustituyeron uno a uno por
% \parencite{clave} en los archivos de secciones/ (mismo resultado
% visual: cita parentética "(Autor, año)").

% ---------- 11. HYPERREF (siempre al final del preámbulo) ----
\usepackage[
  pdfusetitle,
  colorlinks=true,
  linkcolor=black,
  citecolor=black,
  urlcolor=black,
  bookmarksopen=true,
  bookmarksnumbered=true,
  bookmarksdepth=subsection   % marcadores hasta subsección: incluye anexos A-F
]{hyperref}
\usepackage{bookmark}
\urlstyle{same}   % \url y \nolinkurl heredan la fuente del cuerpo, no monoespaciada distinta

% Identificador de variable/script/columna (Instrucción 2.5: "nombres de
% variables con guion bajo deben poder cortarse sin desbordar la columna").
% \nolinkurl (de hyperref) usa el mecanismo de \url --que ya rompe la línea
% en "_", "/", "." y "-"-- pero sin color ni enlace, ideal para nombres de
% columna/variable/script que no son URLs reales.
\newcommand{\var}[1]{\nolinkurl{#1}}

% ---------- 12. UTILIDADES VARIAS ----------------------------
\usepackage{xcolor}
\usepackage{graphicx}
% Diagramas propios del capítulo 7 (flujo CRISP-DM del proyecto),
% dibujados en el propio .tex para que hereden la fuente del cuerpo.
\usepackage{tikz}
\usetikzlibrary{positioning,arrows.meta,calc,fit}
% Columnas de texto alineadas a la izquierda (sin justificar) para
% tablas con celdas estrechas: L{ancho} en lugar de p{ancho} e Y en
% lugar de X de tabularx. Evitan los huecos entre palabras.
\newcolumntype{L}[1]{>{\raggedright\arraybackslash}p{#1}}
\newcolumntype{Y}{>{\raggedright\arraybackslash}X}
% Encabezado de tabla en varias líneas (p. ej. \celda{Bloques\\de prueba})
% para columnas numéricas estrechas sin recurrir a paquetes extra.
\newcommand{\celda}[1]{\begin{tabular}[b]{@{}c@{}}#1\end{tabular}}
% Dato que el autor debe revisar o completar: se imprime en rojo
% oscuro entre corchetes para que no pase inadvertido en el PDF, y
% scripts/pendientes lo lista en docs/PENDIENTES.md.
\newcommand{\pendiente}[1]{%
  \textcolor{red!65!black}{\textbf{[Revisar:}~#1\textbf{]}}}
% Comandos semánticos para reservar espacio en secciones nuevas sin
% escribir formato (reutilizable cuando se transcriba desde Markdown).
\newcommand{\figuraPendiente}[2][\textwidth]{%
  \fbox{\parbox[c][6cm]{#1}{\centering%
    \textcolor{gray!50}{\small Figura pendiente:\\#2}}}%
}
\newcommand{\tablaPendiente}[1]{%
  \fbox{\parbox[c][4cm]{\textwidth}{\centering%
    \textcolor{gray!50}{\small Tabla pendiente:\\#1}}}%
}
\usepackage{enumitem}
\setlist{topsep=4pt,itemsep=2pt}
\frenchspacing
\clubpenalty=10000
\widowpenalty=10000
% Sin \sloppy global: la justificación se resuelve con silabación
% española (polyglossia carga los patrones "spanish") y un margen
% de estiramiento de emergencia acotado, que solo actúa en los
% párrafos que de otro modo desbordarían.
\setlength{\emergencystretch}{1.5em}

% ---------- 13. ENTORNO "RECUADRO" (2.7, definido no usado) ----
% Barra lateral fina + fondo casi blanco, sobrio (estilo "MIT").
% Tres variantes semánticas: hallazgo, decision, advertencia. Se
% define y se documenta aquí, tal como pide la instrucción, pero
% NO se inserta en ningún archivo de secciones/ (eso alteraría el
% texto del autor). El acento de color azul marino del logo
% (Fase 5) todavía no está definido; por ahora la barra usa gris
% neutro, coherente con "el texto siempre negro".
\usepackage{mdframed}
\definecolor{recuadroBarra}{gray}{0.45}
\definecolor{recuadroFondo}{gray}{0.97}
\newmdenv[
  linecolor=recuadroBarra,
  linewidth=2.5pt,
  topline=false, bottomline=false, rightline=false,
  backgroundcolor=recuadroFondo,
  innertopmargin=6pt, innerbottommargin=6pt,
  innerleftmargin=8pt, innerrightmargin=8pt,
  skipabove=8pt, skipbelow=8pt,
]{recuadrobase}
% Uso previsto (no aplicado todavía):
%   \begin{recuadro}{Hallazgo}      ... \end{recuadro}
%   \begin{recuadro}{Decisión}      ... \end{recuadro}
%   \begin{recuadro}{Advertencia}   ... \end{recuadro}
\newenvironment{recuadro}[1]{%
  \begin{recuadrobase}
  \textbf{#1.}\space\ignorespaces
}{%
  \end{recuadrobase}
}
